% !TeX root = ../../main.tex
% =====================================================================
%  Ensayo y adquisición de datos
%  Responsable: 
%  Última revisión: 
% =====================================================================
\chapter{Ensayo y adquisición de datos}
\label{ch:ensayo_adquisicion_datos}

\begin{clave}
\textbf{Capítulo pendiente, y a propósito.} El equipo todavía no tiene campaña de
ensayos ni sistema de adquisición de datos, y este libro no describe métodos que
no practicamos: sería justo el tipo de contenido copiado que hace inútil un
manual interno.

Se redactará cuando haya datos propios que contar. Mientras tanto, los capítulos
que dependen de él ---\ref{ch:piloto_ensayos_pista} (ensayos en pista),
\ref{ch:cfd_validacion} (validación de CFD) y \ref{ch:validacion_termica}
(validación térmica)--- quedan igualmente incompletos, y conviene tenerlo
presente: hoy el equipo diseña con modelos que nadie ha contrastado con la
realidad. Cerrar este capítulo es el paso que convierte el resto del libro en
ingeniería verificada.
\end{clave}

\pendiente{Redactar cuando exista la primera campaña de ensayos. El esqueleto de
apartados de abajo se conserva como guion de partida.}

\section{Qué medir y por qué}
\label{sec:ensayo_adquisicion_datos:01}
\pendiente{Pendiente de redactar.}

\section{Sensores habituales en el coche}
\label{sec:ensayo_adquisicion_datos:02}
\pendiente{Pendiente de redactar.}

\section{Muestreo, filtrado y ruido}
\label{sec:ensayo_adquisicion_datos:03}
\pendiente{Pendiente de redactar.}

\section{Diseño de un ensayo}
\label{sec:ensayo_adquisicion_datos:04}
\pendiente{Pendiente de redactar.}

\section{Del dato a la decisión}
\label{sec:ensayo_adquisicion_datos:05}
\pendiente{Pendiente de redactar.}

% ---------------------------------------------------------------------
%  Cierre del capítulo: los cuatro recuadros son obligatorios.
% ---------------------------------------------------------------------

\begin{caso}
\redactar{Qué hicimos nosotros: decisión tomada, datos de partida, resultado en pista y qué haríamos distinto.}
\end{caso}

\begin{ojo}
\begin{itemize}
  \item \redactar{Error que repetimos cada temporada en este tema.}
\end{itemize}
\end{ojo}

\begin{entregable}
\redactar{Qué produce el lector al terminar: plantilla, hoja de cálculo, checklist o apartado del informe.}
\end{entregable}

\begin{juez}
\begin{enumerate}
  \item \redactar{Pregunta tipo del design event sobre este capítulo.}
\end{enumerate}
\end{juez}
